\honest{The Logical Fallacy of Generalization from Fictional Evidence}{Eliezer Yudkowsky}{2007}

When I bring up advanced AI, ``more than half the time'' people answer with \textsc{The Terminator}, \textsc{The Matrix} or Asimov's robots. I reply that I try to avoid ``the logical fallacy of generalizing from fictional evidence.'' \nb{The post goes on to explain the error by availability, priming, framing and automatic categories, which are psychological effects. The author's chapter on global risks had called it ``this mix of anchoring and availability.''}

But nobody thinks the films are true, you say. My answer is that you can go wrong without believing anything false. First, science fiction writers are storytellers, and ``the requirements of storytelling are not the same as the requirements of forecasting.'' As Nick Bostrom ``points out,'' a film of humanity quietly going extinct would be no fun to watch. \nb{Bostrom offered this ``Good-story bias'' as a supposition that ``could be quite powerful.''} Second, a story is never a rational forecast, because stories do not use probability distributions. I show what a story told in probabilities would sound like: the hero's weapon is a titanium sword with 30\% probability and a crowbar with 20\%. Characters can be ignorant, but the author cannot say ``I don't know.'' The story must follow one line through the future, full of details down to the color of a character's earrings. Then all these details are wrapped up under a short label, which makes them look like a single package.

Third, starting from a film skips the hard part. When the space of possible answers is large, the difficulty is less in checking the right answer than in finding it at all. Asking whether AIs will put us in capsules as in \textsc{The Matrix} jumps to a 100-bit proposition without the 98 bits of evidence needed to make it worth considering. A few more bits after those would make it nearly certain, so nearly all the work is in locating it. That work includes weighing what you do and do not know, widening your confidence intervals, and asking which questions matter.

Starting from a film also frames the debate, and to control the terms of a debate is nearly to control its outcome. \textsc{The Matrix} brings to mind robot armies, ``not a superintelligence snapping nanotechnological fingers.'' \nb{That is also a specific and vivid picture, and the post gives no evidence for it.} It sets up a contest of Us against Them, with questions like ``Who will win?'', in an atmosphere of entertainment. Lost are the many possible designs of mind, the dependence of the future on initial conditions, the power and unpredictability of smarter-than-human intelligence, and people taking the matter seriously. In a gun debate, neither side wants to be introduced as a ``shooting freak'' or a ``victim disarmament advocate''; why accept a frame from Hollywood? Journalists ask whether the future will be like \textsc{2001} or like \textsc{A.I.}, which is as loaded as asking whether to cut veterans' benefits or raise taxes on the rich.

In the ancestral environment, what you saw was true. A glimpse of a single word can prime you, with ``demonstrated strong influence on probability estimates.'' So think what a two-hour film can do. \nb{No study is cited for the single word.} People do not believe the films, ``So far as I can tell'': no one who watched \textsc{The Terminator} hid in a fallout shelter on August 29, 1997. But they recall them as cases that happened somewhere else. ``But didn't that lead to nuclear war in The Terminator?'' is said in the tone of ``didn't that lead to nuclear war on Alpha Centauri?'' The film is not a prophecy but an available historical case. \nb{There was evidence for this that the post does not cite: Green and Brock (2000) found that stories changed readers' beliefs whether they were labeled fact or fiction.}

I caught myself at it once. Someone argued that Vernor Vinge did not expect brain-computer interfaces to raise intelligence much, citing Tunç Blumenthal, the most advanced traveler in \textsc{Marooned in Realtime}, who did not seem very powerful. I replied that Tunç had lost most of his hardware and was crippled. But the question has to be argued in its own right. I could say Vinge chose to depict Tunç as crippled, which gives the choice its proper weight as evidence, but ``There is no was of Tunç Blumenthal.'' My first draft of this post even called \textsc{The Matrix} an ``example.'' A neighboring error is arguing from imaginary evidence: if you did go to the end of the rainbow, you would find a pot of gold, ``which just proves my point!''

Seeing three Borg on screen makes a category, as seeing three zebras does, and the category carries automatic inferences. So people who hear about brain-computer interfaces expect users to be cold and uncompassionate and to walk with heavy mechanical steps. They ask ``Will the future contain Borg?'', not how they know such interfaces make people less nice, or whether they are forming a stereotype on zero evidence. Orwell warned that stock phrases do our thinking for us. The worst damage, ``in my estimation,'' is that other people's imaginings stop us using our own, and I quote Pirsig's student, blocked by repeating what she had heard. \nb{Two days earlier, in ``Original Seeing,'' Yudkowsky quoted the same story and wrote: ``I don't know if this story is based on reality or not, but either way, it's true.''} My last line drops the hedge: remembered fictions are ``the deadliest convenience of all.''
